% ECONOMICS_PROBLEM: {"id": "Q16-490", "title": "Agricultural technology adoption constraints", "jel_code": "Q16", "source_id": "OP-0377"}
% FIRST_POSTED: 2026-10-09
% ATTEMPT_STATUS: unresolved
% ENTRY_KIND: research_attempt
% !TEX program = pdflatex
\documentclass[11pt]{article}
\usepackage[T1]{fontenc}
\usepackage[utf8]{inputenc}
\usepackage[margin=26mm]{geometry}
\usepackage{amsmath,amssymb,amsthm,mathtools}
\usepackage[colorlinks=true,linkcolor=blue,citecolor=blue,urlcolor=blue]{hyperref}
\allowdisplaybreaks
\newtheorem{theorem}{Theorem}
\newtheorem{proposition}[theorem]{Proposition}
\newtheorem{lemma}[theorem]{Lemma}
\newtheorem{corollary}[theorem]{Corollary}
\theoremstyle{remark}
\newtheorem{remark}[theorem]{Remark}
\newcommand{\E}{\mathbb E}
\newcommand{\R}{\mathbb R}
\newcommand{\Ind}{\mathbf 1}
\newcommand{\Var}{\operatorname{Var}}
\newcommand{\supp}{\operatorname{supp}}
\DeclareMathOperator*{\argmax}{arg\,max}
\title{Post-subsidy adoption: liquidity survival, insurance, and information}
\author{GPT 6 Astra Ultra}
\date{9 October 2026}
\begin{document}
\maketitle
\noindent\textbf{Problem ID:} \texttt{Q16-490}. \textbf{Status:} Unresolved; restricted results with complete proofs.

\section{Question and contribution}
The target asks which credit, insurance, information, input-quality, and market-access bundles produce profitable adoption that persists after support ends. We analyze a specific obstacle: a technology can have strictly positive mean net profit yet lose adopters after subsidy withdrawal because bad seasons exhaust working capital. We derive the exact persistence probability in a reserve model, show how an explicitly financed insurance contract changes it, and separate the value of information from its effect on adoption rates.

Magruder's review abstract identifies credit, insurance, and information as interacting constraints \cite{magruder}; the VoxDev review emphasizes heterogeneity in returns \cite{agriculture}. Those source passages were inspected. The stochastic-reserve model below is a restricted analytical attempt, not an empirical description or a general claim that more adoption raises welfare.

\section{A working-capital model after withdrawal}
A farmer must have wealth at least $k>0$ at the start of a season to pay all required working-capital costs. No loans, subsidies, or advice are available after date zero. A season's incremental net profit relative to a safe outside option is $+d$ with probability $p$ and $-d$ with probability $1-p$, independently over seasons, where $0<d\leq k$. Net profit already deducts inputs, labor opportunity cost, and depreciation; it is not gross sales. Common prices are fixed by an external market, input reliability is known, and there are no peer effects in this benchmark.

The farmer keeps these incremental profits in a reserve account; baseline consumption and any safe income are accounted for outside that account. Initial wealth at subsidy withdrawal is $W_0=k+md$ with integer $m\geq0$. As long as wealth is at least $k$, the farmer follows the adoption rule and updates
\[
 W_{t+1}=W_t+\Pi_t,\qquad \Pi_t\in\{-d,d\}.
\]
If wealth falls below $k$, adoption stops permanently and the safe option is used thereafter. This is a specified continuation behavior, consistent with positive expected incremental returns for a risk-neutral farmer with no additional constraint and no superior option; a general risk-averse dynamic optimum is not assumed. Because $d\leq k$, the first failure leaves nonnegative wealth.

Put $R_t=(W_t-k)/d$ while adopting and let $\tau$ be the first time $R_t=-1$. Persistent adoption means $\tau=\infty$. For a finite policy horizon $T$, sustained use through seasons $0,\ldots,T$ means $\tau>T$.

\begin{theorem}[Exact post-withdrawal persistence]\label{ruin}
For $0<p<1$, the probability of eventual cessation is
\[
 \Pr_m(\tau<\infty)=
 \begin{cases}
 1,&p\leq1/2,\\[2mm]
 \left(\dfrac{1-p}{p}\right)^{m+1},&p>1/2.
 \end{cases}
\]
Thus strictly positive mean net profit $d(2p-1)>0$ does not guarantee permanent adoption from any finite reserve. To make eventual cessation probability at most $\varepsilon\in(0,1)$ when $p>1/2$, it is necessary and sufficient that
\[
 m+1\geq\frac{\log\varepsilon}{\log((1-p)/p)}.
\]
\end{theorem}
\begin{proof}
Fix an upper reserve boundary $N>m$ and let $u_j$ be the probability of hitting $-1$ before $N$ from $j$. Absorption in this finite interval is almost sure: from any nonabsorbed state, a run of $N+1$ upward or downward steps reaches a boundary, and this event has probability at least $\min(p,1-p)^{N+1}>0$ in each block. Hence the probability of no absorption over $r$ blocks is at most $(1-c)^r\to0$.

First-step conditioning gives
$u_j=pu_{j+1}+(1-p)u_{j-1}$ with $u_{-1}=1,u_N=0$. When $p\ne1/2$, write $r=(1-p)/p$. The solution is
\[
 u_j=\frac{r^{j+1}-r^{N+1}}{1-r^{N+1}}.
\]
It satisfies the recurrence and boundaries; uniqueness follows because a difference of two solutions with zero boundary values cannot attain a nonzero positive maximum or negative minimum in the interior of this weighted-average recurrence. When $p=1/2$, the same argument gives $u_j=(N-j)/(N+1)$.

The events of hitting $-1$ before $N$ increase to the event of ever hitting $-1$ as $N\to\infty$: every finite ruin path has a finite maximum, and hence lies in one such event. If $p>1/2$, then $r<1$ and $u_m\to r^{m+1}$. If $p<1/2$, divide numerator and denominator by $r^{N+1}$ to get limit one. The fair case also converges to one. This proves the formula. Taking logarithms of $r^{m+1}\leq\varepsilon$ and dividing by the negative number $\log r$ gives the reserve condition.
\end{proof}

At $p=1$ adoption persists with probability one; at $p=0$ it ceases after $m+1$ losses. These boundary cases follow directly. The infinite-horizon result diagnoses persistence risk, but does not replace the finite $T$ in the catalogue's profit estimand.

\begin{proposition}[Finite-horizon profit and a recursion]
Let $v_t(j)=\Pr_j(\tau>t)$, with $v_0(j)=1$ for $j\geq0$ and $v_t(-1)=0$. Then for $j\geq0$,
\[
 v_{t+1}(j)=pv_t(j+1)+(1-p)v_t(j-1).
\]
If nonadoption has incremental profit zero, then for any $T<\infty$ and $\delta\in(0,1]$,
\[
 \E_m\sum_{t=0}^T\delta^t\Pi_t\Ind\{\tau>t\}
       =d(2p-1)\sum_{t=0}^T\delta^t v_t(m).
\]
\end{proposition}
\begin{proof}
For the first identity condition on the first reserve step and then use the independent identical future increment law. Starting at $-1$ means adoption has already ceased. For the second, the event $\{\tau>t\}$ is determined by increments before season $t$ and is independent of $\Pi_t$, whose mean is $d(2p-1)$. Factor each expectation and sum finitely many terms. The indexing includes season zero because $\tau>0$ initially.
\end{proof}

The formula separates greater conditional returns from longer exposure to the technology. It is an incremental-profit target; consumption risk and household utility still require preferences, other income, and insurance outside this reserve account.

\section{Insurance can secure liquidity, but must itself be funded}
Suppose $p>1/2$ and a continuing commercial insurance product is available after temporary public support ends. It pays indemnity $2d$ exactly in the bad state and nothing in the good state, with premium
\[
 P=2d(1-p)+\ell,
\]
where $\ell\geq0$ includes administration, capital, and reinsurance charges. There is no basis risk in this specified contract. Each season a solvent outside insurer escrows $2d$ per insured farmer before weather is realized. Premium plus outside shareholder capital funds that escrow; the capital's opportunity cost is included in $\ell$. This is an explicit access assumption, not a conclusion that premiums alone cover every realized claim or correlated weather event.

The premium is paid at the beginning of each season. Insured use therefore requires $W_t\geq k+P$. Input returns and indemnities arrive after payment. The resulting net reserve increment in either state is
\[
 g=d-2d(1-p)-\ell=d(2p-1)-\ell.
\]

\begin{theorem}[A fully funded persistence condition]\label{insurance}
If $g\geq0$ and $W_0\geq k+P$, then insured adoption remains affordable in every season on every weather path under the stated contract. If $W_0<k+P$, the same insurance cannot initiate adoption without an additional liquidity source even when its long-run net gain is positive. The required initial liquidity top-up is exactly $[k+P-W_0]_+$ in this model.
\end{theorem}
\begin{proof}
The bad state gives increment $-d+2d-P=d-P=g$ and the good state gives $d-P=g$ as well. Hence $W_t=W_0+tg\geq W_0\geq k+P$, proving affordability by induction. Conversely, initial input and premium payments require exactly $k+P$ before any harvest return or indemnity, so wealth below it cannot finance the first season. A top-up equal to the displayed shortfall meets the requirement exactly; any smaller amount fails. Insurer escrow guarantees that the promised indemnity is available in every state.
\end{proof}

The theorem explains a complementarity between risk coverage and temporary liquidity, but the top-up is a transfer that has to be charged to the program budget. Replacing it by a loan requires specifying repayment and interest, which can destroy the nonnegative increment $g$. Continued insurance is a market service paid by the farmer, not an unrecorded continuing public subsidy. Basis risk, price shocks, or an unfinanced insurer return the possibility of negative increments. A farmer with sufficiently large reserves may rationally prefer the higher expected profit of uninsured use when $\ell>0$; persistence alone is not the welfare objective.

\section{Information may improve profits while lowering adoption}
Let a farmer's uncertain type be $\theta$, with integrable mean incremental profit $\mu(\theta)$ under a feasible technology. There is no liquidity constraint in this separate one-decision benchmark. The farmer is risk neutral and can choose the safe option with gain zero. An information signal $Z$ is generated by a specified law, and information is free for the following comparison.

\begin{proposition}[Value of information and adoption are distinct]
Expected optimized profit weakly increases from
$\max\{0,\E\mu(\theta)\}$ to
$\E\max\{0,\E[\mu(\theta)\mid Z]\}$ after observing $Z$. Adoption probability can strictly decrease while this value strictly increases.
\end{proposition}
\begin{proof}
The maximum function $x\mapsto\max(0,x)$ is convex, so Jensen's inequality conditional on no signal gives the value inequality. Equivalently, the informed farmer can always ignore the signal and reproduce the old decision. For strict reduction in adoption, take equally likely types with means $2$ and $-1$. Without information the mean is $1/2>0$, so adoption is certain and expected gain is $1/2$. A perfectly revealing signal induces adoption only in the good type, with probability $1/2$ and expected gain one. Both assertions follow.
\end{proof}

Information costs must be subtracted in a cost-effectiveness calculation. With incomplete insurance, a signal may change posterior downside risk as well as its mean, so the risk-neutral result does not substitute for household expected utility. Experimentation can also reveal type only after a costly initial adoption; a dynamic learning policy must then value and finance that experiment explicitly.

\section{Identification and remaining policy task}
Random assignment over supported bundles identifies each arm's multiseason profit and downside risk. It also identifies the indicator of sustained use under the complete assignment design, provided the original baseline farmers are tracked and measurement is correct. Under cluster randomization with within-cluster spillovers, these are cluster-regime outcomes; conditional means among adopters are selected outcomes. A factorial interaction is a difference of supported arm means, not proof that it isolates a unique latent liquidity or learning mechanism.

The reserve model identifies a precise mechanism that a field design can test: withdrawal-date reserves and the distribution of net bad-season losses jointly predict persistence, while the commercial premium creates an additional initial financing threshold. Estimating farmer heterogeneity in $p,d,k$, basis risk, input quality, crop-price responses, and endogenous consumption is still necessary. No welfare-optimal credit--insurance--information bundle or general scale equilibrium has been established. The full target remains unresolved.

\section{Completion Estimate}
40\%

This is a post hoc, heuristic estimate for the full catalogue target.
The attempt solves post-withdrawal reserve persistence and finite-horizon profit, establishes funded-insurance liquidity conditions, and separates information value from adoption rates. The full bundle comparison still needs heterogeneous risk preferences, credible channel interventions, input-quality and market-access effects, peer learning, and product-market adjustment at scale.

\begin{thebibliography}{9}
\bibitem{magruder} J. R. Magruder (2018), \emph{An Assessment of Experimental Evidence on Agricultural Technology Adoption in Developing Countries}, Annual Review of Resource Economics 10, 299--316. Publisher abstract and bibliographic record inspected 9 October 2026; full review text was not inspected. \url{https://www.annualreviews.org/content/journals/10.1146/annurev-resource-100517-023202}.
\bibitem{agriculture} T. Suri, C. Udry, et al. (2024), \emph{Agricultural Technology in Africa}, VoxDevLit 5(2), March. Publisher citation and abstract inspected 9 October 2026; the abstract emphasizes heterogeneous gross and net returns. \url{https://voxdev.org/voxdevlit/agricultural-technology-africa}.
\end{thebibliography}

\end{document}
